\section*{Resumen}

El fondo cósmico de microondas (CMB) constituye una de las herramientas más importantes para el estudio del Universo primitivo y de la física de la inflación. En particular, la polarización del CMB en el modo-B es sensible a las ondas gravitacionales primordiales (PGWs) generadas durante la época inflacionaria, cuya amplitud está parametrizada por el cociente tensor-escalar $r$. Sin embargo, la detección de este modo-B primordial no depende únicamente de la sensibilidad instrumental, sino que requiere una caracterización precisa de la emisión Galáctica polarizada que lo contamina. Las dos componentes principales que dominan la emisión polarizada del cielo en el rango de microondas son la emisión sincrotrón, producida por electrones relativistas acelerados por el campo magnético Galáctico (GMF) y dominante a bajas frecuencias, y la emisión de polvo térmico, generada por granos de polvo interestelar alineados con el campo magnético y dominante a altas frecuencias. Ambas componentes pueden presentar fracciones de polarización de hasta el $\sim 20\%$ en ciertas regiones del cielo, superando con creces la amplitud de la señal cosmológica. El modelado y caracterización de estos fondos de radiación es por tanto un paso crítico y central en el estudio y diseño de experimentos futuros de polarización del CMB.

\paragraph{}

En este trabajo presentamos una caracterización exhaustiva de las propiedades espectrales de la emisión sincrotrón y del polvo térmico polarizados mediante un análisis del espectro angular de potencias (APS) multifrecuencia. El conjunto de datos utilizado combina los mapas DR1 del instrumento multifrecuencia de QUIJOTE (MFI) a 11 y 17~GHz, las cinco bandas de la misión WMAP (23, 33, 41, 61 y 94~GHz) observadas durante 9 años y las siete bandas de polarización del tercer lanzamiento público de Planck (PR3), que abarcan desde 28.4 hasta 353~GHz. En total, el análisis cubre el rango de frecuencias de 11 a 353~GHz, proporcionando un gran rango frecuencial clave en este análisis en polarización.

\paragraph{}

Los espectros de potencias se calculan en nueve bines de multipolo de anchura $\Delta\ell = 20$ en el rango $20 \leq \ell \leq 200$. Al combinar los autoespectros y los espectros cruzados de las 14 bandas de frecuencia, se obtienen un total de 105 espectros (14 autoespectros y 91 cruzados), lo que genera un conjunto de datos extenso y detallado para el ajuste a un modelo teórico descrito posteriormente.

\paragraph{}

El análisis se realiza sobre el cielo accesible a QUIJOTE, aplicando tres máscaras de análisis construidas a partir de la máscara de análisis de QUIJOTE combinada con cortes en latitud Galáctica de $|b| > 10^\circ$, $|b| > 15^\circ$ y $|b| > 20^\circ$. Cada máscara se apodiza con el esquema C2 de \textsc{NaMaster} con una escala de $5^\circ$, obteniendo fracciones de cielo efectivas de $f_{\rm sky} \approx 0.30$, $0.26$ y $0.23$ respectivamente. El uso de tres máscaras permite evaluar la robustez de los parámetros derivados frente a la contaminación del plano Galáctico.

\paragraph{}

El modelo de ajuste empleado consiste en una descripción simultánea de la emisión sincrotrón y del polvo térmico en el espacio armónico. La componente sincrotrón se modela mediante una ley de potencias en multipolo y en frecuencia, con amplitud $A_s$, índice espectral angular $\alpha_s$ e índice espectral en frecuencia $\beta_s$, este último gobernando la ley de escala de la distribución espectral de energía (SED) en unidades de temperatura de brillo Rayleigh-Jeans, $T_{\rm RJ} \propto \nu^{\beta_s}$. El polvo térmico se modela mediante un cuerpo negro modificado (MBB) con amplitud $A_d$, índice angular $\alpha_d$ e índice espectral $\beta_d$. Un aspecto clave del modelo es la inclusión explícita de la correlación espacial entre ambas componentes, parametrizada mediante el coeficiente $\rho \in [-1,1]$. Los espectros cruzados codifican información complementaria sobre la geometría y coherencia del GMF a lo largo de la línea de visión. La inferencia Bayesiana de los parámetros se realiza mediante un muestreador cadena de Markov Monte Carlo (MCMC), reportando la mediana y los intervalos de credibilidad del 68\% de las distribuciones a posteriori marginalizadas. Se aplican dos estrategias de ajuste complementarias: un ajuste global de ley de potencias que impone una dependencia única sobre todos los bines al multipolo, y un ajuste bin a bin independiente que permite identificar desviaciones dependientes de la escala.

\paragraph{}

Los principales resultados del ajuste de ley de potencias global son los siguientes. Para la componente sincrotrón, el índice espectral en frecuencia es $\beta_s \approx -3.0$ a $-3.1$ para ambos modos de polarización y las tres máscaras, en excelente acuerdo con el valor canónico de la literatura. La inclusión de los datos de QUIJOTE-MFI, que extienden la base frecuencial hasta los 11~GHz, reduce la incertidumbre en $\beta_s$ en un factor de $2.6$-$3.0$ respecto al conjunto WMAP+Planck, siendo el parámetro más beneficiado por la incorporación de QUIJOTE. El cociente B-a-E de la amplitud sincrotrón se obtiene como $r \approx 0.22$-$0.23$, estable para todas las máscaras y en buen acuerdo con estudios previos. Los índices espectrales angulares $\alpha_s$ presentan cierta sensibilidad al rango de multipolos empleado, se muestra que el primer bin ($\ell_{\rm eff} = 29$) es especialmente relevante para determinar el modo BB. Para el polvo térmico, el índice espectral MBB es $\beta_d \approx 1.53$ en todas las configuraciones, en perfecto acuerdo con las estimaciones de Planck, mientras que los índices angulares muestran una asimetría EE-BB significativa ($>5\sigma$). Se caracteriza además la correlación espacial cruzada entre ambas componentes, obteniendo valores pequeños pero significativamente distintos de cero ($\rho \sim 0.1$), con la correlación EE decreciendo con la latitud Galáctica y la BB estable para todas las máscaras, lo que refleja las distintas estructuras del GMF que trazan los dos modos de polarización. El análisis bin a bin corrobora todos los resultados del ajuste global sin evidencia de desviaciones dependientes de la escala angular. La robustez del análisis se verifica adicionalmente con el ajuste utilizando únicamente espectros cruzados. De esta manera se obtienen resultados plenamente consistentes con los obtenidos utilizando además los autoespectros, descartando que los resultados estén condicionados por un sesgo de ruido en los autoespectros.

\paragraph{}

Como trabajo futuro, se contemplan dos extensiones naturales de este análisis. Por un lado, se prevé incorporar la emisión anómala en microondas (AME) al modelo, aprovechando la sensibilidad única de QUIJOTE en el rango de 10-20~GHz para constatar si su contribución polarizada introduce algún sesgo en los parámetros sincrotrón recuperados. Por otro lado, se planea extender el análisis a regiones individuales del cielo de $\sim 5$-$10\%$ de fracción de cielo cada una, con el objetivo de mapear la variabilidad espacial del índice espectral de la emisión sincrotrón $\beta_s$ a nivel de espectro de potencias, proporcionando así restricciones observacionales relevantes para la separación de componentes en futuros experimentos de modo-B.
